\begin{frame}{Maven: знать в лицо}
  \lead{В индустрии часто Maven.} На курсе практическая сборка выполняется
  на Gradle; Maven нужен как ознакомительный промышленный аналог.

  \medskip
  \begin{columns}[T]
    \begin{column}{0.48\textwidth}
      \begin{block}{Что увидите}
        \begin{itemize}
          \item файл \texttt{pom.xml} вместо \texttt{build.gradle.kts};
          \item тот же \texttt{src/main/java};
          \item фазы: \texttt{compile}, \texttt{test}, \texttt{package}.
        \end{itemize}
      \end{block}
    \end{column}
    \begin{column}{0.48\textwidth}
      \begin{block}{Как читать}
        \begin{itemize}
          \item координаты: groupId, artifactId, version;
          \item зависимости --- в XML, не в Kotlin DSL;
          \item IDEA откроет Maven-проект так же, как Gradle.
        \end{itemize}
      \end{block}
    \end{column}
  \end{columns}

  \vspace{0.4em}
  {\footnotesize\color{acGray}
  Практическое задание --- только Gradle. Чужой корпоративный репозиторий
  вы можете встретить уже на Maven.}
\end{frame}
